\noindent Coq development team (2021). Renaming Coq. Announcement, Rocq Prover Discourse forum, 8 April 2021. \href{https://discourse.rocq-prover.org/t/renaming-coq/1264}{Announcement}.
\par\medskip

\noindent Thierry Coquand and Christine Paulin-Mohring (1990). Inductively defined types. In COLOG-88, International Conference on Computer Logic, Tallinn, 1988, Lecture Notes in Computer Science 417, Springer, pp. 50--66. \href{https://doi.org/10.1007/3-540-52335-9_47}{DOI: 10.1007/3-540-52335-9\_47}.
\par\medskip

\noindent Thierry Coquand and G\'erard Huet (1988). The calculus of constructions. Information and Computation 76(2--3), pp. 95--120. \href{https://doi.org/10.1016/0890-5401(88)90005-3}{DOI: 10.1016/0890-5401(88)90005-3}. Huet as a French computer scientist at INRIA: \href{https://en.wikipedia.org/wiki/G\%C3\%A9rard_Huet}{Wikipedia, ``G\'erard Huet''}. Coquand as a French computer scientist and mathematician at the University of Gothenburg: \href{https://en.wikipedia.org/wiki/Thierry_Coquand}{Wikipedia, ``Thierry Coquand''}.
\par\medskip

\noindent Christine Paulin-Mohring. Wikipedia. \href{https://en.wikipedia.org/wiki/Christine_Paulin-Mohring}{Article}. A mathematical logician and computer scientist best known for developing the interactive theorem prover Rocq, formerly Coq.
\par\medskip

\noindent The Rocq Prover. About the Rocq Prover (the name: Inria Rocquencourt and the mythological bird Roc; the old name Coq: the Calculus of Constructions, Thierry Coquand and the rooster). \href{https://rocq-prover.org/about}{Web page}. The project's own answer to why it was called Coq, the French habit of naming research software after animals with Caml, Elan, Foc and Phox as examples, is on the \href{https://github.com/coq/coq/wiki/Presentation}{Presentation page} of its wiki. Coquand's work on the calculus of inductive constructions with Christine Paulin-Mohring: Wikipedia, \href{https://en.wikipedia.org/wiki/Rocq}{Rocq}.
\par\medskip
